% Inputtable proof note. Wrapper: entry-main.tex.
\section{Entry from fixed-target excess}
\label{v5entry:section}

This note isolates the first acquisition step. It does not assert a complete
local exclusion or a sharp Berry--Esseen theorem. Let $P$ be a centered,
unit-variance law, $S_j$ a sum of $j$ independent copies, and set
\[
 \rho=\E|X|^3,\quad \beta=\E[X|X|],\quad \gamma=\E X^3,
 \quad c=\CE=\frac{3+\sqrt{10}}{6\sqrt{2\pi}}.
\]
The threshold $H$ always has the units of the raw sum. A fixed-target
maximizer maximizes
\[
 a(P,H)=\sqrt n\{\Pp(S_n\le H)-\Phi(H/\sqrt n)\}-c\rho.
\]
We use the attainment and contact theorem in the companion fixed-target
note. In particular its closed-support contact
equality, not ratio optimality, is the variational premise here. We consider
a maximizer with $a\ge0$. The actual normalized score is $c+a/\rho$;
no upper bound on it is assumed.

\begin{theorem}[Initial moment and support entry]\label{v5entry:initial}
Assume $n\ge42$, $C_{n-1}\le c$, and the universal structural inequality
\begin{equation}\label{v5entry:structural}
 \sqrt n\,\Delta_n(P)\le\frac{3031}{10000}
                  \left(\rho+\frac{323}{500}\right)
\end{equation}
for standardized laws with finite third moment. Then a fixed-target
maximizer with $a\ge0$ satisfies
\begin{equation}\label{v5entry:entry}
 1\le\rho<\frac{7361}{4000}<\frac{1841}{1000},\qquad
 |\beta|<\frac{321}{500},\qquad
 \operatorname{supp}P\subset(-96/25,96/25).
\end{equation}
The theorem is uniform in all real $H$ and in both signs of a support
point. It requires neither $C_{n-2}$ nor an excess bound of order $1/n$.
\end{theorem}

Inequality \eqref{v5entry:structural} is an external structural input,
recorded with its Shevtsova attribution in the manuscript introduction.
It is not proved by the exact arithmetic below. Only this one of the
two affine structural bounds is needed for the displayed moment radius.

\begin{proof}
Since $c\rho\le\sqrt n\Delta_n(P)$ and $c>819/2000$, elementary division gives
\[
 \rho<\frac{(3031/10000)(323/500)}{819/2000-3031/10000}
       =\frac{7361}{4000}.
\]
The other affine bound $0.3322(\rho+0.429)$ alone would give the larger
number $712569/386500$; it must not be used to claim $7361/4000$.
Put $m=\E X_+=\E X_->0$ and $a_\pm=\E X_\pm^2=(1\pm\beta)/2>0$.
Cauchy on each sign class gives $m\rho\ge a_+^2+a_-^2$ and
$m^2\le a_+a_-$. Hence
\[
 \rho\ge\frac{1+\beta^2}{\sqrt{1-\beta^2}}.
\]
This function increases with $|\beta|$. For $r=1841/1000$, $b=321/500$,
\[
 (1+b^2)^2-r^2(1-b^2)=\frac{93288849}{50000000000}>0,
\]
which proves the second assertion without assuming bounded support.

Write $v=n-1$, $s=\sqrt{n/v}$, $F_v(x)=\Phi(x/\sqrt v)$ and
$e_n=n(s-1)-1/2$. Closed-support contact, the upper predecessor
bound $F_{n-1,P}(H-t)\le F_v(H-t)+c\rho/\sqrt v$, and
$\sqrt n(F_{n,P}(H)-\Phi(H/\sqrt n))=c\rho+a$ give
\begin{equation}\label{v5entry:contact}
 0\le H_n(H,t)+c\{-|t|^3+\tfrac32\rho t^2+3\beta t+\rho e_n\}-na
 \quad(t\in\operatorname{supp}P),
\end{equation}
where exactly
\[
 \frac{H_n(H,t)}{n^{3/2}}
 =F_v(H-t)-F_{v+1}(H)+tF'_{v+1}(H)
                   -\frac{t^2-1}{2}F''_{v+1}(H).
\]
The extra $-na$ can be discarded in this upper bound. This is the precise
reason the predecessor assumption is $C_{n-1}\le c$, not a comparison
with the actual score. No continuity assumption on the distribution
function at $H-t$ is inserted.

Rationalization gives
\[
 ne_n=\frac{s^2(2s+1)}{2(s+1)^2},\qquad
 s<81/80=:S,\qquad
 ne_n<\frac{793881}{2073680}<\frac{48}{125}.
\]
The first expression increases for $s\ge1$, as its logarithmic derivative
is $2/s+2/(2s+1)-2/(s+1)>0$. Thus $e_n<8/875$.

For completeness we rederive the Gaussian estimate used at the radius.
The heat identity $\partial_vF_v=F_v''/2$ permits spatial expansion of
$F_v(H-t)$ to order four, variance expansion of $F_{v+1}$ to order two,
and variance expansion of $F'_{v+1},F''_{v+1}$ to order one. The cubic
coefficient is $(-t^3+3t)/6$, and the quartic coefficient is
$(t^4-6t^2+3)/24$. Therefore, for $t=\pm u$, $z_v=H/\sqrt v$,
\begin{equation}\label{v5entry:heat}
 H_n(H,\pm u)/s^3
 =\pm A(1-z_v^2)\phi(z_v)+B\phi'''(z_v)+\mathcal E,
\end{equation}
where $A=(u^3-3u)/6$, $H_4=u^4-6u^2+3$ and $B=H_4/(24\sqrt v)$.
The four error bounds after division by $s^3$ are respectively
\[
 \frac{u^5\|\phi^{(4)}\|_\infty}{120v},\quad
 \frac{\|\phi^{(5)}\|_\infty}{48v^{3/2}},\quad
 \frac{u\|\phi^{(4)}\|_\infty}{8v},\quad
 \frac{(u^2-1)\|\phi^{(5)}\|_\infty}{16v^{3/2}}.
\]
Gaussian Fourier inversion gives norms at most $3\phi(0)<6/5$ and
$8/\pi<51/20$. Since $v\ge41$ and $\sqrt v>6$, for $u\ge5/2$ this gives
\[
 |\mathcal E|<E(u):=u^5/4100+3u/820+17(3u^2-2)/78720.
\]
These errors are uniform in $H$; they use no higher moment of $P$.

If $A,B>0$ and $B\le A/4$, the leading part of \eqref{v5entry:heat}
is at most $\phi(0)(A+9B^2/(4A))$ for both signs. Indeed, for the
positive sign and $|z_v|\le1$, it is at most
$\phi(0)(A-Az_v^2+3B|z_v|)$, and complete the square. Outside that
interval the first term is nonpositive. For the negative sign the first
term is at most $2e^{-3/2}A\phi(0)<A\phi(0)/2$. In both cases use
$\|\phi'''\|_\infty<2\phi(0)$: on $[0,\sqrt3]$ the polynomial
$3x-x^3$ is at most two, and beyond $\sqrt3$ use
$x^3e^{-x^2/2}\le3\sqrt3e^{-3/2}<2$.

For frozen $r,b$ above define
\begin{align*}
 Q(u)&=-u^3+\tfrac32ru^2+3bu+(8/875)r,\\
 P(u)&=\tfrac25 S^3\left(A^2+\frac{H_4^2}{256\cdot41}\right)
                         +S^3E(u)A+\frac{819}{2000}Q(u)A.
\end{align*}
At $R=96/25<6$ we have $A>0$, $0<H_4<36A$, hence $B<A/4$.
The sign $Q(R)<0$ is checked before replacing $c$ by its lower endpoint.
The upper bound in \eqref{v5entry:contact}, multiplied by $A$, is then
strictly below $P(R)$, and exact arithmetic gives
\[
 -P(R)=\frac{2214258283536204102723}{10250000000000000000000}>\frac15.
\]
To cover each entire exterior ray, differentiation of the exact Gaussian
expression gives
\[
 |\partial_tH_n(H,t)|<\tfrac25(t^2/2+1/2+e_n).
\]
For verification put $y=(H-t)/\sqrt n$, $z=H/\sqrt n$; the derivative is
$n\phi(z)+Ht\phi(z)-ns\phi(sy)$. The second-order Taylor error between
$\phi(y)$ and $\phi(z)$ costs $\phi(0)t^2/2$, while
$\partial_s[s\phi(sy)]=-\phi''(sy)$ costs $\phi(0)n(s-1)$.
Define $L(u)=\tfrac25(u^2/2+1/2+8/875)+(819/2000)Q'(u)$.
At $R$, the six numbers $Q,Q',Q'',L,L',L''$ are strictly negative.
Since $Q'''=-6$ and $L''=2/5-6(819/2000)<0$, all their required signs
persist for $u\ge R$. Thus both radial upper bounds decrease from a
strictly negative boundary value. The contact inequality excludes every
support point with $|t|\ge R$.
\end{proof}

\section{Moments and an actual high-frequency event}
\label{v5entry:moments}
Higher moments are used only after Theorem~\ref{v5entry:initial}.
For $x=|X|$ the nonnegative measure
$d\mu_0=(x-1)^2(x+1/2)\,dP$ has mass $\rho-1$. Direct expansion gives
\begin{align*}
 M_4-1&=\int\frac{(x+1)^2}{x+1/2}\,d\mu_0,\\
 M_5-1&=\int\frac{x^3+2x^2+3x+3/2}{x+1/2}\,d\mu_0,
\end{align*}
where $M_j=\E|X|^j$. Both ratios increase for $x\ge0$: they equal
$x+3/2+1/[4(x+1/2)]$ and
$x^2+3x/2+9/4+3/[8(x+1/2)]$, whose derivatives are nonnegative.
Writing their values at $R=96/25$ as $k_4(R),k_5(R)$, we obtain
\begin{equation}\label{v5entry:mom}
 M_4\le\min(R\rho,1+k_4(R)(\rho-1)),\quad
 M_5\le\min(RM_4,1+k_5(R)(\rho-1)),\quad M_6\le RM_5.
\end{equation}
Also $\gamma^2\le M_4-1$, by expanding
$\E(X^2-1-\gamma X)^2\ge0$.

There is no need to assume a central support gap before the first Fourier
bound. For independent copies let $E=\{(X-X')^2\ge9/25\}$ and let
$M$ be any valid fourth-moment ceiling. Then
\begin{equation}\label{v5entry:event}
 \Pp(E)\ge\alpha:=\frac{(2-9/25)^2}{2M+6}>0,\quad
 |X-X'|\ge3/5\text{ on }E,\quad
 \E[(X-X')^2\mid E]\le2/\alpha.
\end{equation}
Indeed $\E[(X-X')^2;E]\ge2-9/25$ and
$\E(X-X')^4=2M_4+6\le2M+6$, so Cauchy proves the mass bound.
The conditional second-moment bound follows from its unconditional value
two. This is an event under one actual iid marginal, not a replacement
characteristic function.

\section{A sharper first license}
\label{v5entry:license-section}
The retained fixed-target conserved identity is
\begin{equation}\label{v5entry:master}
 (n+H^2)a+nc\rho-\frac{n\phi(H/\sqrt n)}6(\gamma+3J)
                  +n^{3/2}\mathcal L=\mathcal R_c(H).
\end{equation}
Put $\tau(t)=\E[X\mathbf1_{X\ge t}]$ and $\nu=\mathcal L(H-S_{n-1})$.
Over complementary support gaps, define
\[
 J=\sum\tau_{u,v}(v-u)^2,\qquad
 \mathcal L=\sum\tau_{u,v}\int_{[u,v)}(t-u)\,d\nu(t)\ge0.
\]
The half-open interval includes the correct endpoint atoms. The elementary
gap inequality gives
$D=(3+\sqrt{10})\rho-\gamma-3J\ge0$.
In particular the right side of \eqref{v5entry:master} is nonnegative
when $a\ge0$, since it can also be written as
\[
 (n+H^2)a+nc\rho(1-e^{-w})+n\phi(0)e^{-w}D/6+n^{3/2}\mathcal L,
 \qquad w=H^2/(2n).
\]

\begin{lemma}[Moment-cell license]\label{v5entry:license}
Under \eqref{v5entry:master}, with $D,\mathcal L\ge0$, suppose
$n\ge N$, $\rho\in[l,u]$, $a\ge0$, and
$0<c_-\le c$. If a bound valid for every real $H$ gives
\begin{equation}\label{v5entry:gate}
 \mathcal R_c(H)\le T<Nc_-l,
\end{equation}
then either $T<0$ and this cell is excluded, or
\begin{gather}
 \gamma+3J>0,\qquad
 e^{-H^2/(2n)}\ge1-\frac{T}{Nc_-l}>0,\label{v5entry:density}\\
 (n+H^2)a+\frac{n\phi(0)}6D+n^{3/2}\mathcal L
                      \le\mathcal R_c(H)\le T.\label{v5entry:unweighted}
\end{gather}
Thus the first gate need not be the stronger uniform condition $T<2N/5$.
\end{lemma}
\begin{proof}
For a nonexcluded cell $T\ge0$. Divide \eqref{v5entry:master} by $n$:
\[
 \frac{\phi(H/\sqrt n)}6(\gamma+3J)
 =(1+H^2/n)a+c\rho+\sqrt n\mathcal L-\mathcal R_c/n
 \ge c\rho-T/n>0.
\]
This proves the sign. Since $D\ge0$,
$\gamma+3J\le6c\rho/\phi(0)$; hence
$c\rho e^{-w}\ge c\rho-T/n$, proving \eqref{v5entry:density}.
Finally rewrite the conserved expression as the left side of
\eqref{v5entry:unweighted}, plus
$n(1-e^{-w})(c\rho-\phi(0)D/6)$.
Its last coefficient is $\phi(0)(\gamma+3J)/6>0$, so it can be dropped.
\end{proof}

The positivity argument uses no small-defect premise. The input $D\ge0$
is used for the density and unweighted-resource conclusions, not for the
first sign deduction. Thresholds of either sign and $H=0$ are included.
If only $T\le Nc_-l$ is proved, the strict sign and positive density
floor do not follow. The strict gate is a genuine acceptance condition.

\begin{lemma}[First phase price without a small defect]\label{v5entry:phase}
Under Lemma~\ref{v5entry:license}, put
\[
 \mathcal U=(n+H^2)a+n\phi(0)D/6+n^{3/2}\mathcal L,
 \qquad q_{\rm ph}=\tfrac12(c_-l-T/N)>0.
\]
Then $\mathcal U\le\mathcal R_c(H)-q_{\rm ph}H^2$.
In particular a complete polynomial bound
$\mathcal R_c(H)\le A+F|H|-qH^2$ with $F\ge0,q>0$ gives
\[
 \mathcal U\le A+\frac{F^2}{4(q+q_{\rm ph})}.
\]
This improved bound is on $\mathcal U$, not on the raw $\mathcal R_c$.
\end{lemma}
\begin{proof}
Let $M_c=c\rho-\phi(0)D/6$ and $w=H^2/(2n)$.
The exact identity and $\mathcal R_c\le T$ imply
$e^{-w}M_c\ge c\rho-T/n>0$. Therefore
\[
 n(1-e^{-w})M_c=n(e^w-1)e^{-w}M_c
 \ge\frac{H^2}{2}(c\rho-T/n)\ge q_{\rm ph}H^2.
\]
Here $e^w-1\ge w$ for $w\ge0$. Subtract this term from the exact
decomposition $\mathcal R_c=\mathcal U+n(1-e^{-w})M_c$, then complete
the square. There is no premise $D<2$, and no multiplication of separate
lower estimates on $M_c$ and $e^{-w}$.
\end{proof}

\begin{lemma}[Normalized first-cover consequence]\label{v5entry:ten}
Under the preceding hypotheses, suppose a separate complete enclosure
proves $\mathcal R_c(H)/\rho<10$ for all real $H$. For $n\ge42$,
\[
 e^{-H^2/(2n)}>\frac{7199}{17199}>\frac{209}{500},\qquad
 \frac D\rho<\frac{25000}{6979},\qquad
 \frac a\rho<\frac{10}{n}.
\]
In particular the first license is discharged. The actual score satisfies
$c+a/\rho<c+10/n$, but this is an output, not an entry premise.
\end{lemma}
\begin{proof}
Use the pointwise version of the preceding proof with $T=10\rho$:
$10<42(819/2000)<nc$. The density floor is strictly larger than
$1-10/(42c)>1-10/[42(819/2000)]=7199/17199$.
From \eqref{v5entry:unweighted} and $\phi(0)>997/2500$,
$D/\rho<60/(42\phi(0))<25000/6979$.
The nonnegative term $(n+H^2)a$ gives the final bound.
\end{proof}

\section{The first scalar enclosure and its remaining gate}
\label{v5entry:raw}
This section states exactly what a bounded arithmetic cover must supply.
It proves the algebraic reduction, rather than treating an archived PASS
as an analytic premise. Fix a closed cell $[l,u]\subset[1,7361/4000]$.
Choose outward ceilings $M,g,E_5,b$ for $M_4,|\gamma|,M_5,|\beta|$ using
\eqref{v5entry:mom}, $|\gamma|\le\rho$, Pearson, and the sign-moment
inequality. In particular $M<11.1$, $g<1.85$ and $M_6\le RE_5$.

Write $J_1=\sum\tau_{a,b}\int_a^b(t-a)t\,dt$ and
$J_2=\sum\tau_{a,b}\int_a^b(t-a)t|t|\,dt$. Set
\begin{align*}
 K_1&=\E[X|X|^3]-3\beta-3\rho\gamma/2,\\
 K_q&=M_5/4-\rho M_4/2-3\beta\gamma/2+\rho/2,
 \qquad U=J_1+(M_4-3)/6.
\end{align*}
The exact raw remainder in \eqref{v5entry:master} is
\[
 \mathcal R_c=c(3\beta J/2+3\rho J_1-3J_2-K_q)
       +H(cK_1+\phi U)-H^2(c\rho+\gamma\phi/2)+E_n.
\]
The fixed target $c$ is the coefficient in every term. Its enclosing
interval $[819/2000,20487/50000]$ is not a bound on the actual score.

The elementary gap integrals give
$-M_4/6\le J_1\le M_4/2$, $-3J_2-M_5/4\le0$ and $J\le\rho+1$.
The first two follow by integrating, gap by gap,
\[
 \int_a^b(t-a)t\,dt+\tfrac12\int_a^bt^2dt=(b-a)b^2/2,
\]
and its complementary upper inequality; for the negative part of $J_2$
use $\int_u^v(v-x)x^2dx\le\tfrac13\int_u^vx^3dx$.
Consequently a valid constant bound is
\[
 A_2=\tfrac32b(2u+1)+2uM-l/2.
\]
One may take the smaller of $A_2$ and
\[
 A_1=\tfrac32b(2u+1)+\tfrac32\Gamma(R,l)(2u)
                  +\tfrac38u^2(u+1)+\tfrac8{27}u^3-l/2,
\]
where $\Gamma(R,l)$ bounds the nonnegative part of
\[
 -v^2/4+(R+l)v/3-Rl/2+l^3(2R-l)/(12v^2),\qquad l\le v\le R.
\]
This is $v^{-2}\int_l^v(R-x)x(x-l)\,dx$. For each realized negative
magnitude $v$, Tonelli bounds all its gap contributions by this one
integral; multiplying by $v$ and averaging charges $\Gamma\E X_-^3$.
Thus on negative gaps it charges
$3(\rho J_1-J_2)$ by $3\Gamma(\rho-\gamma)/2$; on positive parts use
$t(\rho-t)\le\rho^2/4$, multiply by $(t-a)\ge0$, and integrate
$\sum\tau_{a,b}\int_a^b(t-a)\,dt=J/2$. This gives $3\rho^2J/8$
and includes the positive part of a central gap with $a<0$.
The remaining term uses
$-K_q\le8\rho^3/27+(3b/2-1/2)\rho$, obtained by
$M_5\ge M_4^{3/2}$ and maximizing $\rho y^2/2-y^3/4$.
Thus $A_*:=\min(A_1,A_2)$ is valid when the chosen $\Gamma$ is certified.

Valid bounds on $|K_1|$ are
\begin{align*}
 k_{\rm rad}&=3ub+\frac{R(R+2-3l/2)}{R+1/2}(u-1),\\
 k_{\rm old}&=6u-3u^2/2+3b,\\
 k_R&=Ru/2+\max_{t=l,u}|R/2-3t/2|g+3b.
\end{align*}
For the first, use the radial identity for $K_1+3\rho\beta$ with ratio
$x(x+2-3\rho/2)/(x+1/2)$, convex in $x$ and decreasing in $\rho$.
Its negative part is at most $31/40$, so require its stated endpoint
ceiling to be at least $31/40$. The second uses $R\le6$ and
$|\E[X|X|^3]-3\gamma|\le3\rho$. The third bounds the two fourth
semimoments by $R(\rho\pm\gamma)/2$. Take $k$ to be their minimum.
Also $U\in[-1/2,2M/3-1/2]$, so use
$v_H=(2/5)\max(1/2,2M/3-1/2)$ (or a separately proved smaller bound).
Then $q_*:=c_-l-g/5>0$ is a valid quadratic curvature.

It remains to have a \emph{complete} Fourier bound
\begin{equation}\label{v5entry:fourier-input}
 |E_n|\le K_F+s_F|H|,\quad
 K_F=\tfrac{16}{25}(L_0+K_H+L_2/\sqrt N),\quad s_F=\tfrac{16}{25}L_H.
\end{equation}
The compensated quadratic-hinge inversion has factor $2\sqrt n/\pi$.
The retained triangle transfer uses the parameters $N,M,g,M_6$ and
sixth-order Taylor estimates on $[0,1/2]$. Its conclusion is
\[
 I_0\le(L_0+L_H|H|)/\sqrt n+L_2/n.
\]
The retained complete six-part high-frequency lemma uses the actual
event \eqref{v5entry:event} and gives $\sqrt n I_\infty\le K_H$.
Every damping, floor and window gate in that lemma must hold; one takes
a minimum only of complete six-term bounds. The event proof does not
permit dropping its Gaussian comparator tail. Neither Fourier lemma
uses a score bound or a defect feedback hypothesis.

Completion of the square now proves the all-threshold bound
\begin{equation}\label{v5entry:scalar}
 \mathcal R_c(H)\le
 T:=K_F+c_+A_*+\frac{(c_+k+v_H+s_F)^2}{4q_*},
\end{equation}
provided $A_*\ge0$, $k\ge0$, $q_*>0$ and $c\in[c_-,c_+]$.
This is a scalar inequality, so using a formal coefficient interval
containing $c$ is sound independently of the actual score.
The first numerical obligation is a finite closed cover of
$[1,7361/4000]$ on which the stated Fourier gates and
$T<Nc_-l$ all hold. The endpoint $\rho=1$ is included; at that point
the law is a symmetric two-point law, and all moment ceilings remain
well-defined. Adjacent cell endpoints must both be covered.

\paragraph{Scope of this note.}
The bounded companion checker verifies the support rational reductions,
five exact finite-law moment tests, and a negative control rejecting
radius three with the same envelope. It does not evaluate a whole raw
cover or prove any local comparison predicate. A separate 24-cell
calculation reports $\mathcal R_c/\rho<10$; its source, outward arithmetic,
coverage and complete Fourier premises require their own review before
composition with Lemma~\ref{v5entry:ten}. Subsequent support
contractions, central-gap acquisition, confinement and LOW/HIGH exclusion
still require their own ordered hypotheses. In particular the present
entry theorem alone does not remove the independent finite base.
